\documentclass[../../main.tex]{subfiles}
\begin{document}

{\bf [解答]}

与えられた方程式
\begin{align}
  x^3+px^2+qx+r = 0 \label{eq:1}
\end{align}
について考える．

$e(\theta) = \cos\theta + i\sin\theta$ とする．1以外の2根を $\alpha, \beta$ とすると，
\begin{align}
  \alpha = e(\theta_1), \quad \beta = e(\theta_2) \qquad (0 \le \theta_1 \le \theta_2 < 2\pi) \label{eq:2}
\end{align}
とおける．解と係数の関係から，
\begin{align}
  \begin{cases}
    p = -(1 + e(\theta_1) + e(\theta_2))  \\
    q = e(\theta_1) + e(\theta_2) + e(\theta_1)e(\theta_2) \\
    r = -e(\theta_1)e(\theta_2) = -e(\theta_1 + \theta_2)
  \end{cases} \label{eq:3}
\end{align}
である．$p,q,r$が実数だから，$\alpha,\beta$はどちらも実数か，どちらも複素数かのいずれかである．
それぞれのケースを個別に考える．

\begin{figure}[htb]
  \centering
  \begin{tikzpicture}[scale=1.2, >=stealth, every node/.style={font=\small}]
    % --------------------------------------------------------
    % (左) 実数解の場合: alpha, beta \in R
    % --------------------------------------------------------
    \begin{scope}[shift={(0, 0)}]
      % 座標軸
      \draw[->, thick, gray!80!black] (-1.6, 0) -- (1.8, 0) node[right] {$\mathrm{Re}$};
      \draw[->, thick, gray!80!black] (0, -1.5) -- (0, 1.5) node[above] {$\mathrm{Im}$};
      \node[below left=2pt] at (0, 0) {$O$};

      % 単位円
      \draw[line width=1.1pt, gray!70] (0, 0) circle (1.2);
      \node[above left=1pt, gray!70, font=\footnotesize] at (120:1.2) {$|z|=1$};

      % 根の位置 (±1)
      \node[circle, fill=teal!75!black, inner sep=1.8pt] (R1) at (1.2, 0) {};
      \node[circle, fill=blue!85!black, inner sep=1.8pt] (Rm1) at (-1.2, 0) {};

      \node[above right=1pt, teal!75!black] at (R1) {$1$};
      \node[above left=1pt, blue!85!black] at (Rm1) {$-1$};

      \node[below=4pt, align=center, font=\footnotesize] at (0, -1.6)
      {実数解のとき\\$\alpha, \beta \in \{1, -1\}$};
    \end{scope}

    % --------------------------------------------------------
    % (右) 複素数解（共役）の場合: beta = \bar{alpha}
    % --------------------------------------------------------
    \begin{scope}[shift={(5.2, 0)}]
      % 座標軸
      \draw[->, thick, gray!80!black] (-1.6, 0) -- (1.8, 0) node[right] {$\mathrm{Re}$};
      \draw[->, thick, gray!80!black] (0, -1.5) -- (0, 1.5) node[above] {$\mathrm{Im}$};
      \node[below left=2pt] at (0, 0) {$O$};

      % 単位円
      \draw[line width=1.1pt, gray!70] (0, 0) circle (1.2);
      \node[above left=1pt, gray!70, font=\footnotesize] at (120:1.2) {$|z|=1$};

      % 根 1
      \node[circle, fill=teal!75!black, inner sep=1.8pt] at (1.2, 0) {};
      \node[above right=1pt, teal!75!black] at (1.2, 0) {$1$};

      % 根 alpha, beta (共役, 例として角 55 度)
      \coordinate (A) at (55:1.2);
      \coordinate (B) at (-55:1.2);

      \draw[dashed, blue!60] (A) -- (B);
      \draw[thick, blue!85!black] (0, 0) -- (A);
      \draw[thick, blue!85!black] (0, 0) -- (B);

      \draw[gray!70] (0.4, 0) arc (0:55:0.4);
      \node[font=\footnotesize] at (27.5:0.58) {$\theta_1$};
      \draw[gray!70] (0.4, 0) arc (0:-55:0.4);
      \node[font=\footnotesize] at (-27.5:0.58) {$-\theta_1$};

      \node[circle, fill=blue!85!black, inner sep=1.8pt] at (A) {};
      \node[circle, fill=blue!85!black, inner sep=1.8pt] at (B) {};

      \node[above right=1pt, blue!85!black] at (A) {$\alpha = e^{i\theta_1}$};
      \node[below right=1pt, blue!85!black] at (B) {$\beta = \bar{\alpha}$};

      \node[below=4pt, align=center, font=\footnotesize] at (0, -1.6)
      {虚数解のとき\\$\beta = \bar{\alpha} = e^{-i\theta_1}$};
    \end{scope}
  \end{tikzpicture}
  \caption{単位円上の根 $\alpha, \beta$ の配置（左：実数解の場合，右：共役複素数解の場合）．}
  \label{fig:1973_2_roots_unit_circle}
\end{figure}

\subparagraph{$1^\circ$ $\alpha, \beta \in \mathbb{R}$ の時}

この時，$\alpha,\beta$がとりうる値は$\pm 1$のみだから，あり得る組み合わせは$(\alpha, \beta) = (1, 1), (1, -1), (-1, -1)$ で，\cref{eq:3}より
\begin{align}
  (p, q, r) = (1, -1, -1), (-1, -1, 1), (-3, 3, -1) \label{eq:4}
\end{align}
である．

\subparagraph{$2^\circ$ $\alpha, \beta \notin \mathbb{R}$ の時}

$p, q, r \in \mathbb{R}$ から $\beta = \bar{\alpha}$ となり，$0 \le \theta_1 \le \theta_2 < 2\pi$ とあわせて，$\theta_2 = 2\pi - \theta_1$, $0 < \theta_1 < \pi$ と置ける．
この時\cref{eq:3}より
\begin{align}
  \begin{cases}
    r = -e(2\pi) = -1 \\
    p = -(1 + 2\cos\theta_1) \\
    q = 1 + 2\cos\theta_1
  \end{cases}
\end{align}
だから，
\begin{align}
  p = -q, \quad r = -1, \quad -1 \le q \le 3 \label{eq:5}
\end{align}
である．
\casesend

以上\cref{eq:4,eq:5}をまとめると，\cref{eq:4}の場合の$(p,q,r)=(1,-1,-1),(-3,3,-1)$は\cref{eq:5}に含まれるから，
\begin{align}
  (p = -q, \ r = -1, \ -1 \le q \le 3) \quad \text{or} \quad (p, q, r) = (-1, -1, 1)
\end{align}
である．逆にこの時一根が1で他の二根はその絶対値がいずれも1になるから十分である．

\newpage

\end{document}
